% ============================================================
% 第 1 章  电力系统的基本概念（教材：何仰赞、温增银《电力系统分析》上册 第 1 章）
% 结构对齐教材 1.1—1.5；例题、习题不录
% ============================================================

\chapter{电力系统的基本概念}

电力系统分析研究的是“电能在生产、输送、分配、消费全过程中各运行变量之间的定量关系”。
本章先把对象、单位制（额定值）和网络结构定下来，后面各章的等值电路、潮流、短路、稳定，
都是在这套对象和约定之上展开的。

\section{电力系统的组成}

\subsection{三个层次：动力系统、电力系统、电力网}

\begin{description}[leftmargin=4.5em,style=nextline]
  \item[动力系统] 电力系统加上发电厂的动力部分（锅炉、汽轮机、水轮机、核反应堆、燃气轮机，
        以及燃料输送、除灰、供水等辅助系统）。
  \item[电力系统] 由发电厂、变电所、输电线路、配电线路和用电设备连接而成，完成电能生产、
        输送、分配与消费的统一整体。
  \item[电力网] 电力系统中除去发电机和用电设备以后的部分，即变压器与各级电压线路组成的网络，
        其任务是输送和分配电能。
\end{description}

三者是包含关系：
\[
  \text{动力系统}\;\supset\;\text{电力系统}\;\supset\;\text{电力网}.
\]

\subsection{电能的四个环节}

电能从一次能源到用户，依次经过“发—输—配—用”四个环节：

\begin{enumerate}
  \item \textbf{发电}：火电、水电、核电以及风电、光伏等新能源电源，把一次能源转换为电能。
        发电机端电压一般为 6～20~kV，需经升压变压器升至电网电压。
  \item \textbf{输电}：由最高一级或两级电压的线路构成，把大容量电源的电能送到负荷中心，
        我国常见为 220~kV、330~kV、500~kV、750~kV（交流）及直流输电。
  \item \textbf{配电}：把电能分配到各用户，电压一般为 0.4～35~kV，
        负荷密度大的城市也用 110~kV、220~kV 配电。
  \item \textbf{用电}：异步电动机、同步电动机、电热装置、整流装置、照明设备等。
\end{enumerate}

\begin{table}[htbp]
  \centering
  \caption{电力系统各环节与主要设备}
  \small
  \begin{tabular}{@{}llll@{}}
    \toprule
    环节 & 主要设备 & 典型电压 & 分析中的处理方式 \\
    \midrule
    发电 & 同步发电机、励磁系统 & 6～20~kV & 电势源 $+$ 电抗（第 3 章） \\
    升压/降压 & 双绕组、三绕组、自耦变压器 & 各级 & 阻抗 $+$ 理想变压器（第 2 章） \\
    输电 & 架空线路、电缆 & 110～1000~kV & $\pi$ 型等值电路（第 2 章） \\
    配电 & 架空线路、电缆、配电变压器 & 0.4～35~kV & 简化等值电路 \\
    用电 & 异步电动机为主的综合负荷 & 0.38~kV 等 & 恒定功率或阻抗型负荷（第 9 章） \\
    \bottomrule
  \end{tabular}
\end{table}

\subsection{电力系统运行的三个特点}

\begin{enumerate}
  \item \textbf{电能不能大量存储}。电能的生产、输送、分配和消费几乎同时进行，
        任一时刻发电机发出的有功功率必须等于该时刻负荷所需功率与网络损耗之和。
        这一“功率实时平衡”的约束，是潮流计算、频率调整和稳定问题的共同起点。
  \item \textbf{暂态过程非常短促}。电磁暂态以毫秒计（短路后 $5\sim 10$~ms 内电流即达冲击值），
        机电暂态以秒计。故分析中必须区分电磁暂态与机电暂态，采用不同的数学模型。
  \item \textbf{与国民经济各部门关系极为密切}。供电突然中断后果严重，
        这决定了可靠性要求高于单纯的经济性考虑。
\end{enumerate}

\corepoint{“电能不可存储”是全书的第一性约束：负荷变化必然引起频率与电压变化，
而频率靠有功平衡、电压靠无功平衡来维持。凡是讨论“平衡”的章节（第 11、12、13 章），
根子都在这一条特点上。}

\section{电力系统的额定电压和额定频率}

\subsection{额定电压的定义与约定}

额定电压是电气设备设计和制造的基本技术参数，也是衡量电能质量的基本指标之一。
我国规定，\textbf{电力网的额定电压}就是线路首末两端用电设备所要求的电压，
以 $U_N$ 表示。各类设备的额定电压必须与电网额定电压相适应，但又不完全相同，
原因是电压沿线路有降落、变压器内部有阻抗压降。

\begin{table}[htbp]
  \centering
  \caption{各类电气设备额定电压的确定}
  \small
  \begin{tabular}{@{}p{0.17\linewidth}p{0.34\linewidth}p{0.36\linewidth}@{}}
    \toprule
    设备 & 额定电压的取法 & 理由 \\
    \midrule
    用电设备 & $U_N$（等于所在电网额定电压） & 用电设备直接接于电网，按电网电压设计 \\
    发电机 & $1.05\,U_N$ & 发电机接在线路首端，需补偿线路上的电压降落 \\
    变压器一次绕组 & 等于所接电网的额定电压；\newline 直接与发电机相连时取 $1.05\,U_N$ & 一次绕组相当于用电设备 \\
    变压器二次绕组 & $1.1\,U_N$（短路电压较小或线路很短时可取 $1.05\,U_N$） & 二次额定电压指\textbf{空载}电压；带额定负荷时内部压降约 $5\%$，故空载须高 $10\%$ \\
    \bottomrule
  \end{tabular}
\end{table}

\keyformula{\Delta U\% = \frac{U-U_N}{U_N}\times 100\%,\qquad
            \delta f\% = \frac{f-f_N}{f_N}\times 100\%}

变压器高压绕组常设有分接头，主抽头为 $U_N$，其余为 $\pm 5\%$ 或 $\pm 2\times 2.5\%$
（即 $\pm 5\%$、$\pm 2.5\%$ 若干档），用以在运行中改变变比调压（第 12 章）。
\textbf{实际变比}应取运行时两侧实际抽头的空载线电压之比，而不是铭牌额定电压之比。

\subsection{电压等级与平均额定电压}

我国采用的交流电压等级（线电压有效值）为
\[
  3,\;6,\;10,\;35,\;60,\;110,\;220,\;330,\;500,\;750\ \text{kV},
\]
其中 330~kV 主要用于西北电网，750~kV 用于超高压外送。电压等级按 $\sqrt{3}$ 倍左右的
几何级数配置，使相邻两级之间的输送能力匹配。

在标幺制短路计算中（第 6 章）需要统一基准电压。工程上取各级电网的\textbf{平均额定电压}
$U_{av}\approx 1.05U_N$ 作为基准电压，我国规定为
\begin{equation}
  U_{av}=3.15,\;6.3,\;10.5,\;15.75,\;37,\;115,\;230,\;345,\;525,\;1050\ \text{kV}.
\end{equation}
取平均额定电压作基准后，变压器的额定电压比与基准电压比恰好相等，
变压器的标幺变比等于 $1$，等值电路中理想变压器可以消去——这是简化计算的依据。

\subsection{额定频率与电能质量指标}

我国电力系统的额定频率为 $50$~Hz（工频）。正常运行时允许的频率偏移为
$\pm 0.2\sim \pm 0.5$~Hz（大系统取严、小系统取宽）。供电电压的允许偏移为：
\begin{itemize}
  \item 35~kV 及以上电压级：$\pm 5\%$；
  \item 10~kV 及以下电压级：$\pm 7\%$；
  \item 电压正弦波形畸变率（各次谐波有效值平方和的方根值与基波有效值之比）：
        6～10~kV 供电电压不超过 $4\%$，0.38~kV 不超过 $5\%$。
\end{itemize}

\begin{table}[htbp]
  \centering
  \caption{电能质量的主要指标}
  \small
  \begin{tabular}{@{}lll@{}}
    \toprule
    指标 & 允许值 & 主要影响因素 \\
    \midrule
    频率偏移 & $\pm 0.2\sim\pm 0.5$~Hz & 系统有功功率平衡 \\
    电压偏移 & 35~kV 及以上 $\pm 5\%$；10~kV 及以下 $\pm 7\%$ & 系统无功功率平衡与网络压降 \\
    波形畸变率 & 6～10~kV $\le 4\%$；0.38~kV $\le 5\%$ & 非线性负荷（整流、变频、电弧） \\
    三相不平衡度 & 按国标规定 & 单相负荷、不对称故障 \\
    \bottomrule
  \end{tabular}
\end{table}

\corepoint{频率取决于\textbf{有功}功率平衡，电压取决于\textbf{无功}功率平衡。
系统有功不足则频率偏低，无功不足则电压偏低。因此“保证电能质量”这件事，
在分析上就落实为两个功率平衡问题加两套调整手段（调频：第 13 章；调压：第 12 章）。}

\section{对电力系统运行的基本要求}

对电力系统运行的基本要求可概括为四条：
\begin{enumerate}
  \item 保证安全可靠地供电；
  \item 有合乎要求的电能质量；
  \item 有良好的经济性；
  \item 尽可能减小对生态环境的有害影响。
\end{enumerate}

\subsection{可靠性要求与负荷分级}

按用户对供电可靠性的不同要求，我国将负荷分为三级：

\begin{table}[htbp]
  \centering
  \caption{负荷分级及对供电的要求}
  \small
  \begin{tabular}{@{}p{0.10\linewidth}p{0.45\linewidth}p{0.35\linewidth}@{}}
    \toprule
    级别 & 中断供电的后果 & 供电要求 \\
    \midrule
    第一级 & 危及人身安全；关键设备难以修复的损坏；生产秩序长期不能恢复；
             市政生活重要部门混乱 & 由两个独立电源供电，任一电源故障时另一电源应能保证供电 \\
    第二级 & 大量减产；城市大量居民正常活动受影响 & 一般由双回路供电，条件许可时设备用电源 \\
    第三级 & 不属于第一、二级，短时停电不造成重大损失 & 可由单回路无备用网络供电 \\
    \bottomrule
  \end{tabular}
\end{table}

提高可靠性的措施是“两手”：一是配备足够的有功、无功电源备用，完善网络结构以提高
抵抗扰动的能力（即提高稳定性）；二是严密监视与维修设备、提高运行水平，
并借助计算机进行安全监视与控制。

\subsection{经济性要求}

经济性集中在三个损耗率指标上：
\begin{itemize}
  \item 发电厂的煤耗率（水耗率）——反映发电环节的效率；
  \item 厂用电率——反映发电厂自身用电所占比重；
  \item 电力网的损耗率（网损率）——反映输配电环节的效率。
\end{itemize}
降低网损的措施（提高电压等级、升高运行电压、无功就地补偿、合理选择导线截面、
经济调度等）将在第 14 章系统讨论。

\subsection{电能质量与功率平衡的关系}

电压和频率是电气设备设计与制造的基本参数，也是衡量电能质量的两个基本指标。
它们超差不仅造成废品和减产，还会影响设备安全，严重时危及整个系统。
而要保证电能质量，关键在于“功率平衡 + 配置得当 + 有调整手段”：
\begin{itemize}
  \item 有功功率在额定频率下平衡 $\Rightarrow$ 频率合格；
  \item 无功功率在额定电压下平衡 $\Rightarrow$ 电压合格；
  \item 对谐波污染源进行限制和治理 $\Rightarrow$ 波形合格。
\end{itemize}

\section{电力系统的接线方式}

电力系统的接线包括发电厂主接线、变电所主接线和电力网接线。本节只讨论\textbf{电力网}的
接线方式，它按供电可靠性分为无备用与有备用两大类。

\subsection{无备用接线（开式网络）}

每个负荷只能沿\textbf{唯一路径}取得电能的网络称为\textbf{开式网络}，包括：
\begin{itemize}
  \item \textbf{放射式}：电源点经一条线路直接引至一个负荷点，故障影响范围小、控制方便；
  \item \textbf{干线式}：一条干线沿途分支供电，节省线路走廊；
  \item \textbf{树状}：多级分支，结构最简单。
\end{itemize}
共同优点：简单、设备费用少、运行方便。共同缺点：可靠性低，任一段线路故障或检修
都要中断部分用户供电；干线式和树状网络线路较长时，末端电压往往偏低。

\subsection{有备用接线（闭式网络）}

\begin{itemize}
  \item \textbf{双回路放射式、双回路干线式}：在无备用网络的每一段线路上改用双回路。
        简单、运行方便，可靠性和电压质量明显提高，缺点是设备费用增加很多。
  \item \textbf{环形网络}：由一个或几个电源点与一个或几个负荷点经线路连成环。
        可靠性令人满意且比较经济，缺点是运行调度复杂；单电源环网中一段线路故障开环时，
        正常线段可能过负荷，负荷点电压明显降低。
  \item \textbf{两端供电网络}：可靠性相当于有两个电源的环形网络。
\end{itemize}

环形网络和两端供电网络中，每个负荷点至少经两条线路从不同方向取得电能，
这类网络统称\textbf{闭式网络}。

\begin{table}[htbp]
  \centering
  \caption{两类接线方式的比较}
  \small
  \begin{tabular}{@{}p{0.14\linewidth}p{0.22\linewidth}p{0.22\linewidth}p{0.28\linewidth}@{}}
    \toprule
    接线方式 & 可靠性 & 经济性 & 适用场合 \\
    \midrule
    无备用（放射、干线、树状） & 低，故障即停电 & 最好，设备最省 & 第三级负荷 \\
    双回路 & 较高 & 较差，线路投资大 & 第一、二级负荷比重大的用户 \\
    环网 & 高 & 较好 & 负荷中心的大电厂与枢纽变电所互联 \\
    两端供电 & 高（相当于双电源环网） & 较好 & 有两个独立电源的负荷区 \\
    \bottomrule
  \end{tabular}
\end{table}

\subsection{输电网络与配电网络}

\begin{itemize}
  \item \textbf{输电网络}：任务是把大容量发电厂的电能可靠而经济地送到负荷集中地区，
        通常由系统中最高的一级或两级电压线路组成，连接区域发电厂（经升压站）与枢纽变电所。
        要求：足够的可靠性、满足稳定要求、便于经济调度、对运行方式变更和系统发展有适应性。
        远离负荷中心的大电厂输电干线、向缺乏电源地区供电的干线常采用双回路或多回路；
        位于负荷中心的大电厂和枢纽变电所一般通过环网互联。
  \item \textbf{配电网络}：任务是分配电能，电压一般为 0.4～35~kV，
        大城市负荷密度大时也用 110~kV、220~kV。电源点是相应电压等级的母线，
        负荷点是低一级变电所或用电设备。
\end{itemize}

选择接线方式时须综合考虑：满足用户对可靠性和电压质量的要求、运行灵活方便、
良好的经济指标；一般要对多个可行方案做技术经济比较后确定。

表~\ref{tab:capacity} 给出各级电压单回架空线路输送容量与输送距离的适宜范围，
它反映了“电压等级要与系统规模（容量和供电范围）相适应”这一基本原则。

\begin{table}[htbp]
  \centering
  \caption{各级电压架空线路的输送能力（教材表 1-2）}
  \label{tab:capacity}
  \small
  \begin{tabular}{@{}cccc@{}}
    \toprule
    额定电压/kV & 输送容量/(MV$\cdot$A) & 额定电压/kV & 输送容量/(MV$\cdot$A) \\
    \midrule
    \multicolumn{2}{c}{（输送距离/km）} & \multicolumn{2}{c}{（输送距离/km）} \\
    3   & $0.1\sim1.0$\quad（$1\sim3$）      & 110 & $10\sim50$\quad（$50\sim150$） \\
    6   & $0.1\sim1.2$\quad（$4\sim15$）     & 220 & $100\sim500$\quad（$100\sim300$） \\
    10  & $0.2\sim2$\quad（$6\sim20$）       & 330 & $200\sim800$\quad（$200\sim600$） \\
    35  & $2\sim10$\quad（$20\sim50$）       & 500 & $1000\sim1500$\quad（$150\sim850$） \\
    60  & $3.5\sim30$\quad（$30\sim100$）    & 1000 & $3500\sim5000$\quad（$1000$ 以上） \\
    \bottomrule
  \end{tabular}
\end{table}

\corepoint{无备用接线只允许向第三级负荷供电；第一、二级负荷占较大比重的用户必须由
有备用网络供电。答题时若问“某接线方式是否可行”，先判负荷级别，再判是否满足
$N-1$（任一元件退出仍能供电）。}

\section{电力系统分析课程的主要内容}

\subsection{运行状态与运行变量}

当系统中各种发电、变电、输配电及用电设备之间的联接情况已经确定时，
系统的运行状态由一组\textbf{运行变量}（运行参数）描述：功率、频率、电压、电流、
磁链、电动势以及发电机转子间的相对位移角等。

\begin{table}[htbp]
  \centering
  \caption{稳态与暂态的划分}
  \small
  \begin{tabular}{@{}lll@{}}
    \toprule
    状态 & 含义 & 典型分析内容 \\
    \midrule
    正常稳态 & 三相对称运行，参数在恒定平均值附近微小变化 & 潮流、无功/有功平衡、经济运行 \\
    故障稳态 & 短路进入稳态后的状态（对称或不对称） & 短路电流实用计算 \\
    故障后稳态 & 故障切除后的新稳态 & 稳定分析中的初始点 \\
    暂态 & 从一种运行状态到另一种状态的过渡 & 短路暂态过程、暂态稳定 \\
    振荡 & 运行参数周期性变化的运行状态 & 低频振荡、静态稳定 \\
    \bottomrule
  \end{tabular}
\end{table}

严格说，由于负荷等随机扰动，绝对稳态并不存在；工程上把运行参数只在某恒定平均值附近
作微小变化的状态视为稳态。

\subsection{数学模型：运行变量与系统参数}

分析运行状态的基本方法是：把待研究的状态用数学方程式描述出来，再用适当的数学方法
和计算工具求解。方程式中的未知量是运行变量，系数是\textbf{系统参数}。

系统参数指各元件或其组合在运行中反映其物理特性的参数：电阻、电感（电抗）、
电容（电纳）、时间常数、变压器变比、输入阻抗、转移阻抗等。参数主要取决于元件结构，
也与额定参数密切相关。

以 $n$ 个节点的网络为例，其节点方程为
\keyformula{\dot I_i=\sum_{j=1}^{n}Y_{ij}\dot U_j,\qquad i=1,2,\dots,n}
其中 $Y_{ij}$ 为节点导纳矩阵元素（系统参数），$\dot U_j$ 为节点电压（运行变量）。
这一个式子同时是第 4 章（网络数学模型）、第 6 章（短路电流计算）和第 11 章（潮流计算）
的出发点，区别只在于已知量与未知量的选择不同。

\subsection{上、下册的两条主线}

\begin{table}[htbp]
  \centering
  \caption{课程内容的两条主线}
  \small
  \begin{tabular}{@{}p{0.10\linewidth}p{0.22\linewidth}p{0.28\linewidth}p{0.30\linewidth}@{}}
    \toprule
    分册 & 基本物理量 & 元件处理方式 & 核心问题 \\
    \midrule
    上册（第 1～8 章） & 电流 & 元件当作电路元件，用等值电路代替 & 短路电流计算：各序电抗的物理意义与计算、标幺参数归算、
                                           输入阻抗与转移阻抗、节点导纳/阻抗矩阵 \\
    下册（第 9～19 章） & 功率 & 发电与用电设备从功率转换角度看特性；
                          电网元件着重电压降落与功率损耗 & 稳态：潮流（核心）、无功与电压、有功与频率、经济运行；
                                                                     稳定：功角、静态与暂态稳定、提高稳定的措施 \\
    \bottomrule
  \end{tabular}
\end{table}

上册的关键在\textbf{阻抗（工程计算中主要是电抗）的计算}；下册的关键在\textbf{功率的分布}。
两者共用同一套网络方程，只是把“电流激励—电压响应”换成了“功率注入—电压分布”。

\corepoint{记课程地图：第 2 章给参数、第 3 章给发电机方程、第 4 章给网络方程，
这三章合起来是“建模”；第 5～8 章是“短路”（上册主线）；第 9～14 章是“稳态运行”
（潮流为核）；第 15～19 章是“稳定”。考试大题的骨架通常就是“建模 $\to$ 计算 $\to$ 判据”。}
